\section*{अध्याय \ref{ch:g11:titration} --- अनुमापन}

\begin{solution}{exo:g11:titration:1}
\omterm{def:g11:titration:titration}{अनुमापक} सटीक रूप से ज्ञात सांद्रता वाला विलयन है, जो ब्यूरेट में होता है;
\omterm{def:g11:titration:titration}{अनुमापित विलयन} में मापी जाने वाली स्पीशीज़ होती है, जो फ़्लास्क में होता है।
\omterm{def:g11:titration:equivalence}{तुल्यता} पर मिलाया गया \omterm{def:g11:titration:titration}{अनुमापक} और अनुमापित स्पीशीज़
समीकरण के अनुपात में होते हैं: दोनों खप जाते हैं।
\end{solution}

\begin{solution}{exo:g11:titration:2}
$n = 0.0200 \times 0.0125 = \qty{2.50e-4}{mol}$।
\end{solution}

\begin{solution}{exo:g11:titration:3}
$V_E = \qty{14.3}{mL}$; आयरन(II) की रेखा
\qty{14.3e-4}{mol} से शुरू होती है, अर्थात \qty{1.43e-3}{mol}।
\end{solution}

\begin{solution}{exo:g11:titration:4}
\omterm{def:g11:titration:titration}{अनुमापक} स्वयं रंगीन (बैंगनी) है और उसका उत्पाद लगभग
रंगहीन: परमैंगनेट के अधिकता में होते ही फ़्लास्क में रंग
प्रकट हो जाता है।
\end{solution}

\begin{solution}{exo:g11:titration:5}
5: समीकरण प्रति परमैंगनेट \omterm{def:g9:ions:ion}{आयन} 5 आयरन(II) \omterm{def:g9:ions:ion}{आयन} खपाता है, और
\omterm{def:g11:titration:equivalence}{तुल्यता} पर \omterm{def:g7:chemical-reactions:reactant}{अभिकारक} इन्हीं अनुपातों में
मिलाए जा चुके होते हैं।
\end{solution}

\begin{solution}{exo:g11:titration:6}
$n(\ce{MnO4^-}) = 0.0200 \times 0.00840 = \qty{1.68e-4}{mol}$;
$n(\ce{Fe^{2+}}) = 5 \times \num{1.68e-4} = \qty{8.40e-4}{mol}$;
$c = \num{8.40e-4} / 0.0100 = \qty{0.0840}{mol/L}$;
$C_m = 0.0840 \times 55.8 = \qty{4.69}{g/L}$।
\end{solution}

\begin{solution}{exo:g11:titration:7}
अनुपात 1:1: \qty{10.0}{mL} में $n = 0.0100 \times 0.0114 = \qty{1.14e-4}{mol}$,
इसलिए \qty{100.0}{mL} में \qty{1.14e-3}{mol}, अर्थात
$\num{1.14e-3} \times 176.0 = \qty{0.201}{g}$: लगभग \qty{200}{mg}।
\end{solution}

\begin{solution}{exo:g11:titration:8}
हर बूँद तुरंत अभिक्रिया नहीं करती: \omterm{def:g11:titration:equivalence}{तुल्यता} से पहले \omterm{def:g11:titration:titration}{अनुमापक} का रंग
टिका रहता, और अंत बहुत जल्दी दिखाई देता (या हर बूँद के बाद
मिनटों इंतज़ार करना पड़ता)।
\end{solution}

\begin{solution}{exo:g11:titration:9}
बहुत बड़ा: वह \omterm{def:g11:titration:equivalence}{तुल्यता} से आगे निकल गई है, परमैंगनेट अधिकता में मिलाते हुए।
\omterm{def:g11:titration:equivalence}{तुल्यता} के निकट उसे \omterm{def:g11:titration:titration}{अनुमापक} बूँद-बूँद करके, फ़्लास्क घुमाते हुए,
मिलाना चाहिए था, और पहली स्थायी हल्की गुलाबी आभा पर रुकना चाहिए था।
\end{solution}

\begin{solution}{exo:g11:titration:10}
\qty{10.0}{mL} में $n(\ce{Fe^{2+}}) = 5 \times 0.0200 \times 0.0130 = \qty{1.30e-3}{mol}$:
तनु विलयन के लिए $c = \qty{0.130}{mol/L}$, और
सांद्र विलयन के लिए $10 \times 0.130 = \qty{1.30}{mol/L}$।
\end{solution}

\begin{solution}{exo:g11:titration:11}
$0.0630 \times 55.8 = \qty{3.52}{g}$।
\end{solution}

\begin{solution}{exo:g11:titration:12}
$0.1 / 7.2 \approx \qty{1.4}{\percent}$, जबकि \qty{14.3}{mL} के लिए
\qty{0.7}{\percent}। विलयन के अधिक आयतन (या अधिक
सांद्र नमूने) का \omterm{def:g11:titration:titration}{अनुमापन} कीजिए, या अधिक तनु \omterm{def:g11:titration:titration}{अनुमापक} इस्तेमाल कीजिए।
\end{solution}

\begin{solution}{exo:g11:titration:13}
$n(\ce{MnO4^-}) = 0.0200 \times 0.0176 = \qty{3.52e-4}{mol}$;
\qty{10.0}{mL} में $n(\ce{H2O2}) = \frac{5}{2} \times \num{3.52e-4} = \qty{8.80e-4}{mol}$:
तनु विलयन में \qty{0.0880}{mol/L},
रोगाणुनाशक में \qty{0.880}{mol/L}, अर्थात
$0.880 \times 34.0 = \qty{29.9}{g/L}$ (लगभग 3 \% का विलयन)।
\end{solution}

\begin{solution}{exo:g11:titration:14}
$n(\ce{MnO4^-}) = \num{1.4e-3} / 5 = \qty{2.8e-4}{mol}$, जिसे
लगभग \qty{0.015}{L} में डालना है: $c \approx \qty{0.019}{mol/L}$, इसलिए
\qty{0.0200}{mol/L} का विलयन। दस गुना अधिक सांद्र होने पर $V_E$
लगभग \qty{1.4}{mL} होता: ब्यूरेट की त्रुटि लगभग
\qty{7}{\percent} तक पहुँच जाती।
\end{solution}

\begin{solution}{exo:g11:titration:15}
हाइड्रोजन \omterm{def:g9:ions:ion}{आयनों} के $8 \times 0.0200 \times 0.0143 = \qty{2.29e-3}{mol}$।
\qty{1.0}{mol/L} पर \qty{10}{mL} में \qty{1.0e-2}{mol} होते हैं: पर्याप्त, चार गुने से
भी अधिक। अम्ल के बिना अध्याय की अभिक्रिया लिखे अनुसार
नहीं हो सकती (हाइड्रोजन \omterm{def:g9:ions:ion}{आयन} एक \omterm{def:g7:chemical-reactions:reactant}{अभिकारक} हैं), और परमैंगनेट
अलग ढंग से अभिक्रिया करता।
\end{solution}

\begin{solution}{pb:g11:titration:1}
\textbf{1.} \ce{MnO4^-}/\ce{Mn^{2+}} और \ce{Fe^{3+}}/\ce{Fe^{2+}}।
\omterm{def:g11:redox:oxidant}{ऑक्सीकारक} परमैंगनेट \omterm{def:g9:ions:ion}{आयन} है, \omterm{def:g11:redox:oxidant}{अपचायक} आयरन(II) \omterm{def:g9:ions:ion}{आयन}।

\textbf{2.} \ce{MnO4^- + 8H+ + 5e- -> Mn^{2+} + 4H2O};
\ce{Fe^{2+} -> Fe^{3+} + e-}।

\textbf{3.} दूसरे को 5 से गुणा किया जाता है:
\ce{MnO4^- + 8H+ + 5Fe^{2+} -> Mn^{2+} + 5Fe^{3+} + 4H2O}। बाईं ओर
आवेश $-1 + 8 + 10 = +17$; दाईं ओर $+2 + 15 = +17$।

\textbf{4.} हाइड्रोजन \omterm{def:g9:ions:ion}{आयन} एक \omterm{def:g7:chemical-reactions:reactant}{अभिकारक} हैं: अभिक्रिया के लिए \omterm{def:g9:acids-bases-ph:acidic}{अम्लीय
विलयन} चाहिए।

\textbf{5.} वह पूर्ण है, तेज़ है और फ़्लास्क में परमैंगनेट की एकमात्र
अभिक्रिया है; बैंगनी परमैंगनेट \omterm{def:g11:titration:equivalence}{तुल्यता} से पहले खप जाता है और
उसके बाद फ़्लास्क को गुलाबी रंग देता है।

\textbf{6.} ब्यूरेट में; नवचंद्रक के निचले भाग पर, आँख की
ऊँचाई पर, पहले और बाद में, \qty{0.05}{mL} तक पढ़ा जाता है।

\textbf{7.} हल्का पीला (आयरन \omterm{def:g9:ions:ion}{आयन}); हर बैंगनी बूँद अभी भी मौजूद
आयरन(II) \omterm{def:g9:ions:ion}{आयनों} द्वारा तुरंत खप जाती है।

\textbf{8.} $0.0200 \times 0.0143 = \qty{2.86e-4}{mol}$।

\textbf{9.} $n(\ce{Fe^{2+}}) = 5 \times \num{2.86e-4} =
\qty{1.43e-3}{mol}$।

\textbf{10.} $\num{1.43e-3} \times 55.8 = \qty{0.0798}{g} =
\qty{79.8}{mg}$।

\textbf{11.} $(80 - 79.8)/80 \approx \qty{0.3}{\percent}$।

\textbf{12.} \qty{14.3}{mL} में से $\qty{0.1}{mL}$ \qty{0.7}{\percent} है,
लगभग \qty{0.6}{mg}: परिणाम, $79.8 \pm 0.6$ mg, लेबल से
मेल खाता है।

\textbf{13.} $14.1$: $5 \times 0.0200 \times 0.0141 \times 55.8 =
\qty{78.7}{mg}$; $14.4$: \qty{80.4}{mg}।

\textbf{14.} $(79.8 + 78.7 + 80.4)/3 = \qty{79.6}{mg}$।

\textbf{15.} गोलियाँ स्वयं अपनी आयरन मात्रा में थोड़ी भिन्न हैं,
लगभग \qty{1}{\percent} तक।

\textbf{16.} नहीं: $V_E$ लगभग \qty{2.9}{mL} होता, और ब्यूरेट की
त्रुटि \qty{3.5}{\percent}, पाँच गुना बदतर।

\textbf{17.} विटामिन C भी परमैंगनेट द्वारा ऑक्सीकृत होता: अभिक्रिया
अब एकमात्र नहीं रहती, $V_E$ बहुत बड़ा होता, और
आयरन का अधिक आकलन होता।

\textbf{18.} $\num{1.43e-3} \times 151.9 = \qty{0.217}{g}$, लगभग
\qty{217}{mg}।

\textbf{19.} $\num{1.43e-3} \times \num{6.02e23} = \num{8.6e20}$ \omterm{def:g9:ions:ion}{आयन}।

\textbf{20.} प्रति गोली लगभग \qty{80}{mg} आयरन: लेबल
ईमानदार है।
\end{solution}
